\section{Συγκριτική Ανάλυση \ENG{E-Learning} \ENG{Standards}}
\label{sec:comparative_analysis}

Η τελική σύγκριση του κεφαλαίου οργανώνεται γύρω από την αρχιτεκτονική, τα λειτουργικά χαρακτηριστικά, τα σενάρια χρήσης, το κόστος και την τεχνική πολυπλοκότητα. Η παρούσα ενότητα δίνει έμφαση σε συγκριτικούς πίνακες και διαγράμματα, ώστε η αντιπαραβολή των βασικών προτύπων να είναι άμεση και πυκνή σε πληροφορία.

\subsection{Πολυδιάστατη Σύγκριση}
\label{subsec:multidimensional_comparison}

Το Σχήμα~\ref{fig:standards_positioning} και ο Πίνακας~\ref{tab:architecture_comparison} αποτυπώνουν τη θέση κάθε προτύπου ως προς τη σύζευξη με το \ENG{LMS}, την ευελιξία του μοντέλου δεδομένων και τον βασικό μηχανισμό ανταλλαγής πληροφορίας.

\begin{figure}[H]
  \centering
  \begin{chapterfigurebox}
\begin{tikzpicture}[
  scale=0.92,
  every node/.style={font=\small},
  axislabel/.style={font=\small\bfseries},
  helper/.style={font=\footnotesize, text=black!70},
  standard/.style={rectangle, rounded corners, draw, thick, minimum width=2.2cm, minimum height=0.9cm, align=center, fill=white, drop shadow}
]

\draw[->, very thick] (0,0) -- (10.4,0);
\draw[->, very thick] (0,0) -- (0,6.4);
\draw[dashed, gray!50] (5.1,0) -- (5.1,6.1);
\draw[dashed, gray!50] (0,3.05) -- (10.1,3.05);

\node[axislabel] at (5.2,-0.75) {Βαθμός αποσύζευξης από το \ENG{LMS}};
\node[helper] at (1.5,-0.3) {Σφιχτή σύζευξη};
\node[helper] at (8.8,-0.3) {Ανεξάρτητες υπηρεσίες};

\node[axislabel, rotate=90] at (-0.95,3.2) {Ευελιξία μοντέλου δεδομένων};
\node[helper, rotate=90] at (-0.35,1.0) {Σταθερό μοντέλο};
\node[helper, rotate=90] at (-0.35,5.1) {Ευέλικτο μοντέλο};

\node[standard, fill=aiccblue!25] (aicc) at (1.4,1.0) {\ENG{AICC}};
\node[standard, fill=scormgreen!25] (scorm12) at (3.1,1.5) {\ENG{SCORM} 1.2};
\node[standard, fill=scormgreen!45] (scorm2004) at (3.5,2.5) {\ENG{SCORM} 2004};
\node[standard, fill=imslightblue!35] (lti) at (8.1,1.7) {\ENG{LTI}};
\node[standard, fill=accentpurple!35] (cmi5) at (6.8,4.0) {\ENG{cmi5}};
\node[standard, fill=xapiorange!45] (xapi) at (8.5,5.1) {\ENG{xAPI}};

\draw[-Latex, thick, scormgreen] (scorm2004) to[bend left=15] (xapi);
\draw[-Latex, thick, xapiorange] (xapi) to[bend left=15] (cmi5);

\node[helper, align=center, text width=2.7cm] at (2.0,5.2) {Πακέτα περιεχομένου με ισχυρή εξάρτηση από \ENG{LMS/browser}};
\node[helper, align=center, text width=2.7cm] at (8.2,3.2) {Υπηρεσιακή αρχιτεκτονική με έμφαση σε \ENG{events} και ενσωμάτωση};

\end{tikzpicture}
\end{chapterfigurebox}

  \caption{Συγκριτική τοποθέτηση βασικών προτύπων ως προς αρχιτεκτονική σύζευξη και ευελιξία δεδομένων}
  \label{fig:standards_positioning}
\end{figure}

\begin{table}[H]
\centering
\caption{Συγκριτική αποτύπωση αρχιτεκτονικής, μεταφοράς δεδομένων και σύζευξης}
\label{tab:architecture_comparison}
\small
\setlength{\tabcolsep}{5pt}
\renewcommand{\arraystretch}{1.25}
\begin{chaptertablebox}
\begin{tabularx}{\textwidth}{|p{1.8cm}|Y|Y|Y|Y|Y|}
\hline
\textbf{Πρότυπο} & \textbf{Μοντέλο εκτέλεσης} & \textbf{Μεταφορά δεδομένων} & \textbf{Κύριος αποθηκευτής} & \textbf{Ισχυρό σημείο} & \textbf{Βασικός περιορισμός} \\
\hline
\ENG{AICC} & \ENG{Server-linked external units} & \ENG{HACP} \ENG{request/response} & \ENG{LMS/AICC server} & Πρώιμη διαλειτουργικότητα μεταξύ ετερογενών συστημάτων & Παρωχημένο πρωτόκολλο και περιορισμένη ευελιξία μοντέλου \\
\hline
\ENG{SCORM} 1.2 & \ENG{Browser session} μέσα στο \ENG{LMS} & \ENG{JavaScript API} & \ENG{LMS runtime} & Απλή και ώριμη συμβατότητα περιεχομένου & Περιορισμένο \ENG{tracking} και αυστηρή εξάρτηση από \ENG{browser} \\
\hline
\ENG{SCORM} 2004 & \ENG{Browser session} με \ENG{sequencing engine} & \ENG{JavaScript API} + \ENG{sequencing rules} & \ENG{LMS runtime} & Προηγμένο \ENG{sequencing} και κανόνες πλοήγησης & Υψηλή πολυπλοκότητα συγγραφής, δοκιμών και αποσφαλμάτωσης \\
\hline
\ENG{xAPI} & Αποσυζευγμένοι πελάτες δραστηριοτήτων & \ENG{REST/JSON statements} & \ENG{LRS} & Πλούσια καταγραφή εμπειριών σε πολλαπλά περιβάλλοντα & Απαιτεί ισχυρότερο σχεδιασμό \ENG{verbs}, \ENG{activities} και \ENG{analytics} \\
\hline
\ENG{cmi5} & \ENG{Structured launch} πάνω σε \ENG{xAPI} & \ENG{Launch contract} + \ENG{xAPI statements} & \ENG{LRS} & Γεφυρώνει την πειθαρχία του \ENG{SCORM} με την ευελιξία του \ENG{xAPI} & Μικρότερο οικοσύστημα εργαλείων και χαμηλότερη διάχυση \\
\hline
\ENG{LTI} & \ENG{Tool launch} από το \ENG{LMS} & \ENG{OIDC/OAuth/JWT} & \ENG{LMS} + \ENG{external tool} & Ομαλή ενσωμάτωση τρίτων εργαλείων και \ENG{single sign-on} & Δεν αποτελεί πρότυπο περιεχομένου ή εγγενούς \ENG{tracking} \\
\hline
\end{tabularx}
\end{chaptertablebox}
\end{table}


\subsection{Σύγκριση Λειτουργικών Χαρακτηριστικών}
\label{subsec:functional_comparison}

Ο Πίνακας~\ref{tab:standards_functional_comparison} συνοψίζει τα βασικά λειτουργικά γνωρίσματα των κυριότερων προτύπων και αναδεικνύει πού υπερτερεί το καθένα σε επίπεδο \ENG{tracking}, πλοήγησης, ασφάλειας και διαλειτουργικότητας.

\begingroup
\setlength{\LTpre}{0pt}
\setlength{\LTpost}{0pt}
\setlength{\tabcolsep}{4pt}
\renewcommand{\arraystretch}{1.22}
\footnotesize
\begin{longtable}{>{\raggedright\arraybackslash}p{3.7cm} >{\centering\arraybackslash}p{1.3cm} >{\centering\arraybackslash}p{1.3cm} >{\centering\arraybackslash}p{1.3cm} >{\centering\arraybackslash}p{1.3cm} >{\centering\arraybackslash}p{1.3cm} >{\centering\arraybackslash}p{1.3cm}}
\caption{Λειτουργική Σύγκριση \ENG{E-Learning} \ENG{Standards}}
\label{tab:standards_functional_comparison}\\
\toprule
\textbf{Χαρακτηριστικό} & \textbf{\ENG{AICC}} & \textbf{\ENG{SCORM} 1.2} & \textbf{\ENG{SCORM} 2004} & \textbf{\ENG{xAPI}} & \textbf{\ENG{cmi5}} & \textbf{\ENG{LTI}} \\
\midrule
\endfirsthead

\caption[]{Λειτουργική Σύγκριση \ENG{E-Learning} \ENG{Standards} (συνέχεια)}\\
\toprule
\textbf{Χαρακτηριστικό} & \textbf{\ENG{AICC}} & \textbf{\ENG{SCORM} 1.2} & \textbf{\ENG{SCORM} 2004} & \textbf{\ENG{xAPI}} & \textbf{\ENG{cmi5}} & \textbf{\ENG{LTI}} \\
\midrule
\endhead

\midrule
\multicolumn{7}{r}{\footnotesize Συνεχίζεται στην επόμενη σελίδα}\\
\endfoot

\bottomrule
\endlastfoot

\multicolumn{7}{l}{\cellcolor{gray!15}\textit{Αρχιτεκτονική \& Περιβάλλον}} \\
\midrule
\ENG{Browser-based} & Όχι & Ναι & Ναι & Όχι & Όχι & Ναι \\
\ENG{Mobile} \ENG{support} & Περιορ. & Όχι & Όχι & Ναι & Ναι & Ναι \\
\ENG{Offline} \ENG{tracking} & Ναι* & Όχι & Όχι & Ναι & Ναι & Όχι \\
\ENG{Cloud-ready} & Όχι & Όχι & Όχι & Ναι & Ναι & Ναι \\
\midrule
\multicolumn{7}{l}{\cellcolor{gray!15}\textit{Δεδομένα \& Καταγραφή}} \\
\midrule
\ENG{Data} \ENG{model} & Σταθερό & Σταθερό & Σταθερό & Ευέλικτο & Περιορ. & \ENG{N/A} \\
\ENG{Extensions} & Όχι & Όχι & Περιορ. & Ναι & Ναι & Ναι \\
\ENG{Rich} \ENG{context} & Όχι & Όχι & Όχι & Ναι & Ναι & Μέτριο \\
\ENG{Historical} \ENG{data} & Περιορ. & Περιορ. & Περιορ. & Ναι & Ναι & Όχι \\
\midrule
\multicolumn{7}{l}{\cellcolor{gray!15}\textit{\ENG{Sequencing} \& \ENG{Navigation}}} \\
\midrule
\ENG{Simple} \ENG{sequencing} & Ναι & Ναι & Ναι & Όχι* & Ναι & Όχι \\
\ENG{Advanced} \ENG{seq.} & Όχι & Όχι & Ναι & Όχι* & Μέτριο & Όχι \\
\ENG{Adaptive} \ENG{paths} & Όχι & Όχι & Ναι & Όχι* & Όχι & Όχι \\
\midrule
\multicolumn{7}{l}{\cellcolor{gray!15}\textit{Διαλειτουργικότητα}} \\
\midrule
\ENG{LMS} \ENG{integration} & Πολύπλ. & Εύκολο & Μέτριο & Εύκολο & Εύκολο & Πολύ εύκολο \\
\ENG{External} \ENG{tools} & Όχι & Όχι & Όχι & Ναι & Ναι & Ναι \\
\ENG{Cross-platform} & Περιορ. & Περιορ. & Περιορ. & Ναι & Ναι & Ναι \\
\midrule
\multicolumn{7}{l}{\cellcolor{gray!15}\textit{Ασφάλεια}} \\
\midrule
\ENG{Authentication} & Βασικό & Βασικό & Βασικό & \ENG{OAuth} 2.0 & \ENG{OAuth} 2.0 & \ENG{OAuth} 2.0 \\
\ENG{Encryption} & Όχι & Όχι & Όχι & \ENG{HTTPS} & \ENG{HTTPS} & \ENG{HTTPS} \\
\ENG{Privacy} \ENG{controls} & Όχι & Όχι & Όχι & Ναι & Ναι & Ναι \\
\midrule
\multicolumn{7}{l}{\cellcolor{gray!15}\textit{Υιοθέτηση \& Υποστήριξη}} \\
\midrule
\ENG{Market} \ENG{adoption} & Χαμηλή & Υψηλή & Μέτρια & Αυξ. & Χαμηλή & Υψηλή \\
\ENG{Tool} \ENG{support} & Παλαιό & Εξαιρ. & Καλό & Καλό & Αναπτ. & Εξαιρ. \\
\ENG{Documentation} & Περιορ. & Εξαιρ. & Καλό & Εξαιρ. & Καλό & Εξαιρ. \\
\ENG{Active} \ENG{dev.} & Όχι & Όχι & Όχι & Ναι & Ναι & Ναι \\
\end{longtable}
\endgroup



\subsection{Σύγκριση Περιπτώσεων Χρήσης}
\label{subsec:use_case_comparison}

Η επιλογή προτύπου εξαρτάται κυρίως από το είδος της εμπειρίας μάθησης και το επίπεδο ελέγχου που απαιτεί ο οργανισμός. Το Σχήμα~\ref{fig:standard_selection_flow} λειτουργεί ως συνοπτικός οδηγός επιλογής, ενώ ο Πίνακας~\ref{tab:use_case_mapping} αντιστοιχίζει τα συνηθέστερα σενάρια με την προσφορότερη τεχνολογική λύση.

\begin{figure}[H]
  \centering
  \makebox[\textwidth][c]{%
    \resizebox{0.96\textwidth}{!}{%
      \begin{chapterfigurebox}
\begin{tikzpicture}[
  node distance=1.4cm and 1.8cm,
  every node/.style={font=\small},
  decision/.style={diamond, aspect=2.1, draw, thick, align=center, text width=4.2cm, inner sep=4pt, fill=gray!10},
  choice/.style={rectangle, rounded corners, draw, thick, minimum width=2.5cm, minimum height=0.95cm, align=center, fill=white, drop shadow},
  line/.style={-Latex, thick},
  yn/.style={font=\footnotesize, fill=white, inner sep=1pt}
]

\node[decision] (tools) {Χρειάζεται ενσωμάτωση εξωτερικού εργαλείου στο \ENG{LMS};};
\node[choice, fill=imslightblue!35, right=4.8cm of tools] (lti) {\ENG{LTI}};

\node[decision, below=of tools] (package) {Χρειάζεται πακεταρισμένο μάθημα με κλασική συμβατότητα \ENG{LMS};};
\node[decision, below left=1.6cm and 2.3cm of package] (sequencing) {Χρειάζεται σύνθετο \ENG{sequencing};};
\node[choice, fill=scormgreen!25, below left=1.5cm and 1.8cm of sequencing] (scorm12) {\ENG{SCORM} 1.2};
\node[choice, fill=scormgreen!45, below right=1.5cm and 1.8cm of sequencing] (scorm2004) {\ENG{SCORM} 2004};

\node[decision, below right=1.6cm and 2.3cm of package] (launch) {Χρειάζεται τυποποιημένο \ENG{launch/completion contract};};
\node[choice, fill=accentpurple!35, below left=1.5cm and 1.8cm of launch] (cmi5) {\ENG{cmi5}};
\node[choice, fill=xapiorange!45, below right=1.5cm and 1.8cm of launch] (xapi) {\ENG{xAPI}};

\draw[line] (tools) -- node[yn, above] {Ναι} (lti);
\draw[line] (tools) -- node[yn, right] {Όχι} (package);

\draw[line] (package) -- node[yn, above left] {Ναι} (sequencing);
\draw[line] (package) -- node[yn, above right] {Όχι} (launch);

\draw[line] (sequencing) -- node[yn, left] {Όχι} (scorm12);
\draw[line] (sequencing) -- node[yn, right] {Ναι} (scorm2004);

\draw[line] (launch) -- node[yn, left] {Ναι} (cmi5);
\draw[line] (launch) -- node[yn, right] {Όχι} (xapi);

\node[font=\footnotesize, align=center, text width=10.5cm, below=1.5cm of package] {
Συνήθης πορεία μετάβασης: \ENG{SCORM} 1.2 $\rightarrow$ \ENG{cmi5} $\rightarrow$ \ENG{xAPI}, ενώ το \ENG{LTI} λειτουργεί συμπληρωματικά στο επίπεδο ενσωμάτωσης εργαλείων.
};

\end{tikzpicture}
\end{chapterfigurebox}
%
    }%
  }
  \caption{Συνοπτική ροή επιλογής προτύπου με βάση το εκπαιδευτικό και τεχνικό ζητούμενο}
  \label{fig:standard_selection_flow}
\end{figure}

\begin{table}[H]
\centering
\caption{Αντιστοίχιση τυπικών σεναρίων χρήσης με τα καταλληλότερα πρότυπα}
\label{tab:use_case_mapping}
\small
\setlength{\tabcolsep}{5pt}
\renewcommand{\arraystretch}{1.25}
\begin{chaptertablebox}
\begin{tabularx}{\textwidth}{|p{3.4cm}|p{2.0cm}|p{2.0cm}|Y|}
\hline
\textbf{Σενάριο} & \textbf{Κύρια επιλογή} & \textbf{Εναλλακτική} & \textbf{Αιτιολόγηση} \\
\hline
\ENG{Compliance} εκπαίδευση με τυπικά \ENG{modules} & \ENG{SCORM} 1.2 & \ENG{cmi5} & Προτεραιότητα έχει η καθολική συμβατότητα με \ENG{LMS}, η χαμηλή πολυπλοκότητα και η ώριμη υποστήριξη από \ENG{authoring tools}. \\
\hline
Προσαρμοστικές μαθησιακές διαδρομές με σύνθετο \ENG{sequencing} & \ENG{SCORM} 2004 & \ENG{cmi5} & Όταν απαιτούνται \ENG{branching}, \ENG{rollup rules} και προαπαιτούμενα, το \ENG{SCORM} 2004 παραμένει το πιο πλήρες τυπικό μοντέλο. \\
\hline
\ENG{Cross-platform} και \ENG{offline tracking} & \ENG{xAPI} & \ENG{cmi5} & Η ανάγκη για καταγραφή σε \ENG{mobile}, \ENG{simulations} και φυσικά περιβάλλοντα ευνοεί το \ENG{LRS}-κεντρικό μοντέλο του \ENG{xAPI}. \\
\hline
Μετάβαση από υφιστάμενο \ENG{SCORM} χαρτοφυλάκιο & \ENG{cmi5} & \ENG{xAPI} & Το \ENG{cmi5} διατηρεί πειθαρχημένο \ENG{launch/completion contract} και άρα λειτουργεί ως ελεγχόμενη γέφυρα μετάβασης. \\
\hline
Ενσωμάτωση εξωτερικών εργαλείων στο \ENG{LMS} & \ENG{LTI} & -- & Το ζητούμενο εδώ δεν είναι η τυποποίηση περιεχομένου αλλά η ασφαλής εκκίνηση εργαλείων, η μεταφορά ρόλων και η επιστροφή βαθμολογίας. \\
\hline
Ενοποιημένο μαθησιακό ιστορικό από πολλαπλές πηγές & \ENG{xAPI} & \ENG{LTI} + \ENG{xAPI} & Όταν ο οργανισμός θέλει ενιαίο \ENG{learner record}, το \ENG{xAPI} υπερέχει ως κοινό σχήμα συμβάντων και το \ENG{LTI} λειτουργεί επικουρικά για το επίπεδο ενσωμάτωσης. \\
\hline
\end{tabularx}
\end{chaptertablebox}
\end{table}


\subsection{Εξελικτικές Σχέσεις και Συνέχεια}
\label{subsec:evolutionary_relationships}

Η συγκριτική θεώρηση επιβεβαιώνει ότι τα πρότυπα δεν λειτουργούν ως απομονωμένες τεχνολογικές νησίδες, αλλά ως διαδοχικά επίπεδα κληρονομικότητας, εξειδίκευσης και συμπλήρωσης. Η αλληλουχία αυτή αποτυπώνεται στο Σχήμα~\ref{fig:standards_genealogy}.

\begin{figure}[H]
  \centering
  % FIGURE 2.8: Standards Genealogy Tree (Evolution & Relationships)
% Directed graph showing standards evolution from 1980 to 2026
\begin{chapterfigurebox}
\begin{tikzpicture}[
  scale=0.85,
  every node/.style={font=\small},
  standard/.style={rectangle, rounded corners, draw, thick, minimum width=2.2cm, minimum height=1cm, align=center, fill=white, drop shadow},
  influence/.style={-latex, thick, dashed},
  evolution/.style={-latex, very thick},
  year/.style={font=\tiny, gray}
]

% Root: CBT (1980s)
\node[standard, fill=problemred!20] at (0, 8) (cbt) {
  \textbf{\ENG{CBT}}\\
  \tiny \ENG{Computer-Based}\\
  \tiny \ENG{Training}
};
\node[year, above] at (cbt.north) {1980\ENG{s}};

% First generation split: AICC & IMS
\node[standard, fill=aiccblue!30] at (-4, 6) (aicc) {
  \textbf{\ENG{AICC}}\\
  \tiny \ENG{Aviation} \ENG{Industry}\\
  \tiny \ENG{CBT} \ENG{Committee}
};
\node[year, left] at (aicc.west) {1988};

\node[standard, fill=imslightblue!30] at (4, 6) (ims) {
  \textbf{\ENG{IMS}}\\
  \tiny \ENG{Global} \ENG{Learning}\\
  \tiny \ENG{Consortium}
};
\node[year, right] at (ims.east) {1997};

% Evolution arrows from CBT
\draw[evolution, problemred] (cbt) -- (aicc);
\draw[evolution, problemred] (cbt) -- (ims);

% Metadata branch (parallel development)
\node[standard, fill=metadatagray!30, text width=2.5cm] at (0, 5.5) (metadata) {
  \textbf{\ENG{Dublin} \ENG{Core}}\\
  \textbf{\ENG{IEEE} \ENG{LOM}}\\
  \tiny \ENG{Metadata} \ENG{Standards}
};
\node[year, above right] at (metadata.north east) {1995-2002};
\draw[influence, gray] (cbt) -- (metadata);

% SCORM (convergence of AICC + ADL + IMS)
\node[standard, fill=scormgreen!40, text width=2.3cm] at (-2, 3.5) (scorm1) {
  \textbf{\ENG{SCORM} 1.0-1.2}\\
  \tiny \ENG{ADL} \ENG{Initiative}
};
\node[year, left] at (scorm1.west) {2000-2001};

\draw[evolution, aiccblue] (aicc) -- (scorm1);
\draw[influence, gray] (metadata) -- (scorm1);
\draw[influence, imslightblue] (ims) -- (scorm1);

% SCORM maturity
\node[standard, fill=scormgreen!60] at (-2, 1.5) (scorm2004) {
  \textbf{\ENG{SCORM} 2004}\\
  \tiny 4\ENG{th} \ENG{Edition}
};
\node[year, left] at (scorm2004.west) {2004-2009};

\draw[evolution, scormgreen] (scorm1) -- (scorm2004);

% IMS branches: QTI, CC, LTI
\node[standard, fill=imslightblue!40, text width=2cm] at (5, 4) (qti) {
  \textbf{\ENG{IMS} \ENG{QTI}}\\
  \tiny \ENG{Question} \&\\
  \tiny \ENG{Test}
};
\node[year, above] at (qti.north) {2000};

\node[standard, fill=imslightblue!50, text width=2cm] at (5, 2.5) (cc) {
  \textbf{\ENG{IMS} \ENG{CC}}\\
  \tiny \ENG{Common}\\
  \tiny \ENG{Cartridge}
};
\node[year, above] at (cc.north) {2008};

\node[standard, fill=imslightblue!60] at (5, 1) (lti) {
  \textbf{\ENG{IMS} \ENG{LTI}}\\
  \tiny \ENG{Learning} \ENG{Tools}\\
  \tiny \ENG{Interoperability}
};
\node[year, above] at (lti.north) {2010};

\draw[evolution, imslightblue] (ims) -- (qti);
\draw[evolution, imslightblue] (qti) -- (cc);
\draw[evolution, imslightblue] (cc) -- (lti);

% xAPI Revolution (influenced by SCORM limitations)
\node[standard, fill=xapiorange!50, text width=2.2cm] at (0, -0.5) (xapi) {
  \textbf{\ENG{xAPI}}\\
  \tiny (\ENG{Tin} \ENG{Can} \ENG{API})\\
  \tiny \ENG{Experience} \ENG{API}
};
\node[year, left] at (xapi.west) {2011};

\draw[influence, scormgreen, bend left=20] (scorm2004) to node[midway, above, font=\tiny, sloped] {\ENG{limitations}} (xapi);
\draw[influence, aiccblue, bend right=30] (aicc) to (xapi);

% cmi5 (AICC + xAPI convergence)
\node[standard, fill=accentpurple!50, text width=2cm] at (-3, -2) (cmi5) {
  \textbf{\ENG{cmi5}}\\
  \tiny \ENG{AICC} + \ENG{xAPI}
};
\node[year, below] at (cmi5.south) {2016};

\draw[evolution, xapiorange] (xapi) -- (cmi5);
\draw[influence, aiccblue, bend left=45] (aicc) to node[midway, below, font=\tiny, sloped] {\ENG{legacy} \ENG{support}} (cmi5);

% Contemporary convergence: Total Learning Architecture
\node[standard, fill=blue!20, text width=3.5cm, minimum height=1.5cm] at (1.5, -3.5) (tla) {
  \textbf{\ENG{Total} \ENG{Learning} \ENG{Architecture}}\\
  \tiny \ENG{xAPI} 2.0, \ENG{LRS}, \ENG{Analytics},\\
  \tiny \ENG{AI/ML}, \ENG{Open} \ENG{Badges} 3.0
};
\node[year, below] at (tla.south) {2020-2026};

\draw[evolution, xapiorange] (xapi) -- (tla);
\draw[evolution, accentpurple] (cmi5) -- (tla);
\draw[influence, imslightblue, bend left=30] (lti) to (tla);
\draw[influence, scormgreen, dashed] (scorm2004) to node[midway, right, font=\tiny] {\ENG{legacy}} (tla);

% Key relationships legend
\node[below, font=\tiny, align=left] at (0, -5.5) {
  \textbf{\ENG{Legend}:} \\
  \tikz \draw[evolution] (0,0) -- (0.5,0); \ENG{Direct} \ENG{evolution} \quad
  \tikz \draw[influence] (0,0) -- (0.5,0); \ENG{Influence/inspiration} \quad
  \textcolor{problemred}{\textbullet} \ENG{Pre-standard} \quad
  \textcolor{aiccblue}{\textbullet} \ENG{AICC} \ENG{family} \quad
  \textcolor{scormgreen}{\textbullet} \ENG{SCORM} \ENG{family} \quad
  \textcolor{xapiorange}{\textbullet} \ENG{xAPI} \ENG{family} \quad
  \textcolor{imslightblue}{\textbullet} \ENG{IMS} \ENG{family}
};

% \ENG{Timeline} annotations (side)
\node[right, font=\scriptsize, align=left, text width=3cm] at (7.5, 5) {
  \textbf{1980-1995:}\\
  \tiny \ENG{Fragmentation}\\
  \tiny \ENG{era}
};

\node[right, font=\scriptsize, align=left, text width=3cm] at (7.5, 2.5) {
  \textbf{1997-2010:}\\
  \tiny \ENG{Standardization}\\
  \tiny \& \ENG{convergence}
};

\node[right, font=\scriptsize, align=left, text width=3cm] at (7.5, -1) {
  \textbf{2011-2020:}\\
  \tiny \ENG{Disruption}\\
  \tiny (\ENG{xAPI})
};

\node[right, font=\scriptsize, align=left, text width=3cm] at (7.5, -3.5) {
  \textbf{2020+:}\\
  \tiny \ENG{Modern}\\
  \tiny \ENG{integration}
};

\end{tikzpicture}
\end{chapterfigurebox}

  \caption{Γενεαλογικό δέντρο \ENG{e-learning} προτύπων και βασικών γραμμών εξέλιξης}
  \label{fig:standards_genealogy}
\end{figure}

\subsection{Οικονομική και Επιχειρησιακή Διάσταση}
\label{subsec:economic_dimension}

Η επιχειρησιακή αποτίμηση δεν περιορίζεται στις τεχνικές δυνατότητες αλλά επεκτείνεται στο κόστος ανάπτυξης, συντήρησης, εκπαίδευσης ομάδων και εξάρτησης από εργαλεία. Ο Πίνακας~\ref{tab:development_costs} συνοψίζει τη σχετική οικονομική επιβάρυνση ανά πρότυπο.

\begin{table}[H]
\centering
\caption{Συγκριτικό Κόστος Ανάπτυξης ανά \ENG{Standard}}
\label{tab:development_costs}
\small
\setlength{\tabcolsep}{5pt}
\renewcommand{\arraystretch}{1.25}
\begin{chaptertablebox}
\begin{tabularx}{\textwidth}{|Y|>{\centering\arraybackslash}p{2.2cm}|>{\centering\arraybackslash}p{2.2cm}|>{\centering\arraybackslash}p{2.2cm}|>{\centering\arraybackslash}p{2.2cm}|}
\hline
\textbf{\ENG{Standard}} & \textbf{\ENG{Initial} \ENG{Dev}} & \textbf{\ENG{Maintenance}} & \textbf{\ENG{Tools} \ENG{Cost}} & \textbf{\ENG{Training}} \\
\hline
\ENG{SCORM} 1.2 & Χαμηλό & Χαμηλό & Χαμηλό & Χαμηλό \\
\ENG{SCORM} 2004 & Μέτριο & Μέτριο-Υψηλό & Μέτριο & Μέτριο \\
\ENG{xAPI} & Μέτριο-Υψηλό & Μέτριο & Μέτριο & Υψηλό \\
\ENG{cmi5} & Μέτριο & Μέτριο & Μέτριο & Μέτριο \\
\ENG{LTI} & Χαμηλό* & Χαμηλό & Χαμηλό & Χαμηλό \\
\hline
\end{tabularx}
\end{chaptertablebox}
\end{table}



\subsection{Τεχνική Πολυπλοκότητα}
\label{subsec:technical_complexity}

Η διαφορά ανάμεσα σε ένα θεωρητικά πλούσιο και σε ένα πρακτικά υιοθετήσιμο πρότυπο αποτυπώνεται καλύτερα όταν εξεταστούν μαζί η πολυπλοκότητα υλοποίησης, η ωριμότητα εργαλείων και ο ρόλος κάθε προτύπου σε στρατηγικές μετάβασης.

\begin{table}[H]
\centering
\caption{Μήτρα τεχνικής πολυπλοκότητας, ωριμότητας εργαλείων και ρόλου μετάβασης}
\label{tab:complexity_adoption_matrix}
\small
\setlength{\tabcolsep}{5pt}
\renewcommand{\arraystretch}{1.25}
\begin{chaptertablebox}
\begin{tabularx}{\textwidth}{|p{1.8cm}|p{2.1cm}|p{2.2cm}|p{2.0cm}|p{2.4cm}|Y|}
\hline
\textbf{Πρότυπο} & \textbf{Πολυπλοκότητα υλοποίησης} & \textbf{Ωριμότητα εργαλείων} & \textbf{Δυσκολία \ENG{debugging}} & \textbf{Ωριμότητα οικοσυστήματος} & \textbf{Τυπικός ρόλος σε στρατηγική μετάβασης} \\
\hline
\ENG{AICC} & Μέτρια & Χαμηλή & Μέτρια & Χαμηλή & Καθαρά \ENG{legacy} τεχνολογία, χρήσιμη κυρίως για κατανόηση της προέλευσης του \ENG{SCORM}. \\
\hline
\ENG{SCORM} 1.2 & Χαμηλή & Πολύ υψηλή & Χαμηλή & Πολύ υψηλή & Παραμένει η βάση των ιστορικών καταλόγων μαθημάτων και το σημείο εκκίνησης πολλών οργανισμών. \\
\hline
\ENG{SCORM} 2004 & Υψηλή & Μέτρια & Υψηλή & Μέτρια & Επιλέγεται μόνο όταν η ανάγκη για τυποποιημένο \ENG{sequencing} είναι ουσιώδης και σαφώς τεκμηριωμένη. \\
\hline
\ENG{xAPI} & Μέτρια-Υψηλή & Μέτρια & Μέτρια & Αυξανόμενη & Στοχεύει σε σύγχρονα \ENG{analytics}, πολυκαναλική παρακολούθηση και μακροχρόνιο \ENG{learner record}. \\
\hline
\ENG{cmi5} & Μέτρια & Χαμηλή-Μέτρια & Μέτρια & Αναδυόμενη & Λειτουργεί ως πρακτικό μονοπάτι μετάβασης για οργανισμούς που θέλουν \ENG{xAPI} χωρίς να χάσουν πειθαρχημένο \ENG{course launch}. \\
\hline
\ENG{LTI} & Μέτρια & Υψηλή & Μέτρια & Υψηλή & Συμπληρωματικός άξονας \ENG{integration}, όχι άμεσος αντικαταστάτης των προτύπων περιεχομένου. \\
\hline
\end{tabularx}
\end{chaptertablebox}
\end{table}


\subsection{Διδάγματα Σχεδιασμού}
\label{subsec:design_lessons}

Η ιστορική σύγκριση οδηγεί σε επαναλαμβανόμενα σχεδιαστικά συμπεράσματα που αφορούν τόσο την υπερβολική πολυπλοκότητα όσο και την ανάγκη για σαφή διαχωρισμό \ENG{launch}, \ENG{tracking} και \ENG{tool integration}. Τα συμπεράσματα αυτά συνοψίζονται στον Πίνακα~\ref{tab:design_lessons_summary}.

\begin{table}[H]
\centering
\caption{Σύνοψη σχεδιαστικών διδαγμάτων από τη συγκριτική εξέλιξη των προτύπων}
\label{tab:design_lessons_summary}
\small
\setlength{\tabcolsep}{5pt}
\renewcommand{\arraystretch}{1.25}
\begin{chaptertablebox}
\begin{tabular}{|p{3.2cm}|p{6cm}|p{4.8cm}|}
\hline
\textbf{Παρατήρηση} & \textbf{Τεκμήριο στη σύγκριση} & \textbf{Σχεδιαστική συνέπεια} \\
\hline
Απλό \ENG{launch contract} & Το \ENG{SCORM} 1.2 και το \ENG{cmi5} διαδόθηκαν ευκολότερα από προσεγγίσεις που απαιτούν εκτεταμένη παραμετροποίηση. & Η ελάχιστη υποχρεωτική ροή εκκίνησης πρέπει να είναι μικρή, προβλέψιμη και εύκολα ελέγξιμη. \\
\hline
Αποσύζευξη καταγραφής από \ENG{runtime} & Το \ENG{xAPI} υπερέχει όταν η καταγραφή πρέπει να συνεχίζεται πέρα από ένα \ENG{browser session} ή ένα μόνο \ENG{LMS}. & Η σύλληψη συμβάντων, η αποθήκευση και η ανάλυση πρέπει να παραμένουν διακριτά επίπεδα. \\
\hline
\ENG{Tooling} πριν από πολυπλοκότητα & Το \ENG{SCORM} 2004 έχασε δυναμική επειδή το \ENG{sequencing} ξεπέρασε την πραγματική υποστήριξη από εργαλεία συγγραφής και αποσφαλμάτωσης. & Καμία προδιαγραφή δεν πρέπει να απαιτεί περισσότερα από όσα μπορεί να υποστηρίξει το διαθέσιμο οικοσύστημα εργαλείων. \\
\hline
Οι γέφυρες αποδίδουν περισσότερο από τις απότομες ρήξεις & Το \ENG{cmi5} είναι χρήσιμο ακριβώς επειδή συνδυάζει οικεία \ENG{course} συμβόλαια με \ENG{xAPI} καταγραφή. & Οι στρατηγικές μετάβασης πρέπει να επιτρέπουν συνύπαρξη \ENG{legacy} και νέων υλοποιήσεων. \\
\hline
Η ενσωμάτωση εργαλείων είναι ξεχωριστό πρόβλημα από το \ENG{content packaging} & Το \ENG{LTI} επιτυγχάνει υψηλή υιοθέτηση χωρίς να επιχειρεί να αντικαταστήσει \ENG{SCORM}, \ENG{xAPI} ή \ENG{cmi5}. & Το \ENG{launch}, το \ENG{tracking} και η ενσωμάτωση εργαλείων πρέπει να αντιμετωπίζονται ως διακριτά στρώματα διαλειτουργικότητας. \\
\hline
\end{tabular}
\end{chaptertablebox}
\end{table}


\subsection{Συμπεράσματα Συγκριτικής Ανάλυσης}
\label{subsec:comparative_conclusions}

Συνολικά, το \ENG{SCORM} 1.2 παραμένει η πιο πρακτική λύση για κλασικά \ENG{LMS}-κεντρικά μαθήματα, το \ENG{SCORM} 2004 διατηρεί αξία όπου απαιτείται αυστηρό \ENG{sequencing}, το \ENG{xAPI} αποτελεί την πιο ευέλικτη επιλογή για πολυκαναλική καταγραφή εμπειριών και το \ENG{cmi5} προσφέρει ρεαλιστική γέφυρα μετάβασης ανάμεσα σε δομημένο \ENG{course launch} και σύγχρονο \ENG{event tracking}. Το \ENG{LTI}, τέλος, δεν ανταγωνίζεται τα πρότυπα περιεχομένου αλλά συμπληρώνει το οικοσύστημα ως τυποποιημένο στρώμα ενσωμάτωσης εργαλείων.

Το βασικό συμπέρασμα είναι ότι η σύγχρονη διαλειτουργικότητα δεν οικοδομείται γύρω από ένα μοναδικό «καθολικό» πρότυπο, αλλά γύρω από συνδυασμούς προτύπων που καλύπτουν διακριτά επίπεδα της μαθησιακής πλατφόρμας: πακετάρισμα και \ENG{runtime}, καταγραφή εμπειριών, και ολοκλήρωση εξωτερικών υπηρεσιών.
